\documentclass[a4paper, twoside]{article}
\setcounter{secnumdepth}{5}
\setcounter{tocdepth}{5}
\usepackage[english]{babel}
\usepackage{textcomp}
\usepackage{amsmath,amsthm,amsfonts,amssymb,epsfig}
\usepackage{array}
\usepackage{datetime}
\usepackage{lipsum}% http://ctan.org/pkg/lipsum
\usepackage[left=1.1in,top=1in,right=1.1in]{geometry}% http://ctan.org/pkg/geometry
\usepackage{listings}% http://ctan.org/pkg/listings
\usepackage{spverbatim}
\usepackage{hyperref}
\usepackage{microtype}
\hypersetup{colorlinks=true, urlcolor=black, linkcolor=black}
\usepackage{graphicx}
\graphicspath{ {images/} }
\usepackage{parskip}
\usepackage{titlesec} %used for diminishing heading sizes
\usepackage[square, sort, comma, numbers]{natbib} %%uses titles for cited references
\usepackage[fit]{truncate}
\usepackage{fancyhdr}
\pagestyle{fancy}
\fancyhead{}
\fancyhead[RO, RE]{\thepage}
\fancyhead[LO, LE]{\rightmark}
\renewcommand{\sectionmark}[1]{\markboth{}{\textsc{\thesection~#1}}}
\fancyfoot[C]{}%hide footer
\usepackage{xcolor} 
\usepackage[scaled=1]{couriers}

\xdefinecolor{gray}{rgb}{0.6,0.6,0.6} 


\titleformat*{\section}{\LARGE\bfseries\sffamily}
\titleformat*{\subsection}{\Large\bfseries\sffamily}
\titleformat*{\subsubsection}{\large\bfseries\sffamily}
\titleformat*{\paragraph}{\large\bfseries\sffamily}
\titleformat*{\subparagraph}{\large\bfseries\sffamily}
\renewcommand{\familydefault}{\sfdefault} %sans-serif font



\begin{document}


%----------------------------------------------------------------------
% Definition for "lstlisting" blocks
%----------------------------------------------------------------------
% --- USAGE ---
%
% \begin{lstlisting}[style=R}
% ...
% \end{lstlisting}
%
% % \begin{lstlisting}[style=output}
% ...
% \end{lstlisting}
%----------------------------------------------------------------------

% By default, make listings all black so it's easy to spot the ones that aren't set to a style.
% This is just a debugging technique.
\lstset{backgroundcolor=\color{black}}

% Define scala language first
% ``define'' Scala
\lstdefinelanguage{scala}{
  morekeywords={abstract,case,catch,class,def,%
    do,else,extends,false,final,finally,%
    for,if,implicit,import,match,mixin,%
    new,null,object,override,package,%
    private,protected,requires,return,sealed,%
    super,this,throw,trait,true,try,%
    type,val,var,while,with,yield},
  otherkeywords={=>,<-,<\%,<:,>:,\#,@},
  sensitive=true,
  morecomment=[l]{//},
  morecomment=[n]{/*}{*/},
  morestring=[b]``,
  morestring=[b]',
  morestring=[b]''``
}

\lstdefinestyle{R}{
  language=R,
  frame=single,
  breaklines,
  basicstyle=\ttfamily,
  commentstyle=\textbf,% comment style
  keywordstyle=\ttfamily,
  numbers=left,% display line numbers on the left side 
  numberstyle=\scriptsize,% use small line numbers 
  numbersep=10pt,% space between line numbers and code
  backgroundcolor=\color{white}, 
  showstringspaces=false % don't show spaces as weird char.
}

\lstdefinestyle{python}{
  language=python,
  frame=single,
  breaklines,
  basicstyle=\ttfamily,
  commentstyle=\textsl,% comment style
  keywordstyle=\ttfamily,
  numbers=left,% display line numbers on the left side 
  numberstyle=\scriptsize,% use small line numbers 
  numbersep=10pt,% space between line numbers and code
  backgroundcolor=\color{white}, 
  showstringspaces=false %don't show spaces as weird char.
}

\lstdefinestyle{Scala}{
  language=scala,
  frame=single,
  breaklines,
  basicstyle=\ttfamily,
  commentstyle=\textsl,% comment style
  keywordstyle=\ttfamily,
  numbers=left,% display line numbers on the left side 
  numberstyle=\scriptsize,% use small line numbers 
  numbersep=10pt,% space between line numbers and code
  backgroundcolor=\color{white}, 
  showstringspaces=false % don't show spaces as weird char.
}

\lstdefinestyle{Bash}{
  language=bash,
  frame=single,
  breaklines,
  basicstyle=\ttfamily,
  commentstyle=\textsl,% comment style
  keywordstyle=\ttfamily,
  numbers=left,% display line numbers on the left side 
  numberstyle=\scriptsize,% use small line numbers 
  numbersep=10pt,% space between line numbers and code
  backgroundcolor=\color{white}, 
  showstringspaces=false % don't show spaces as weird char.
}


\definecolor{mygray}{rgb}{0.92,0.92,0.92}

\lstdefinestyle{output}{
  frame=single,
  breaklines,
  basicstyle=\ttfamily,
  numbers=left,% display line numbers on the left side 
  numberstyle=\scriptsize,% use small line numbers 
  numbersep=10pt,% space between line numbers and code
  backgroundcolor=\color{mygray}, 
  showstringspaces=false %don't show spaces as weird char.
}

\newcommand{\waterExampleInR} {
\textbf{Example in R} \\
}

\newcommand{\waterExampleInPython} {
\textbf{Example in Python} \\
}


\thispagestyle{empty} %removes page number
\newgeometry{bmargin=0cm, hmargin=0cm}

\begin{center}
\textsc{\Large\bf{Fast Scalable R with H2O}}
\\
\bigskip
\textsc{\small{Patrick Aboyoun \hspace{40pt} Spencer Aiello \hspace{40pt} Anqi Fu \hspace{40pt} Mark Landry \hspace{40pt} Eric Eckstrand \hspace{40pt}  Jessica Lanford}}
\\
\bigskip
\line(1,0){250}  %inserts  horizontal line

{\url{http://h2o.gitbooks.io/r-with-h2o/}}

\bigskip
August 2015: Third Edition 
\\%add front page image here? (wavy lines)
\bigskip
\end{center}

\null\vfill
\begin{figure}[!b]
\noindent\makebox[\textwidth]{%
\centerline{\includegraphics[width=\paperwidth]{waves.png}}}
\end{figure}

\newpage
\restoregeometry

\null\vfill %move next text block to lower left of new page
\thispagestyle{empty} %removes page number


{\raggedright\vfill\ 

Fast Scalable R with H2O\\
  by Patrick Aboyoun,  Spencer Aiello, Anqi Fu, Mark Landry \&\ Jessica Lanford\\
\bigskip
  Published by H2O.ai, Inc. \\
2307 Leghorn St. \\
Mountain View, CA 94043\\
\bigskip
\textcopyright 2015 H2O.ai, Inc. All Rights Reserved. 
\bigskip

August 2015: Third Edition
\bigskip

Photos by \textcopyright H2O.ai, Inc. 
\bigskip

While every precaution has been taken in the\\
preparation of this book, the publisher and\\
authors assume no responsibility for errors or\\
omissions, or for damages resulting from the\\
use of the information contained herein.\\
\bigskip
Printed in the United States of America. 


}\par

\newpage
\tableofcontents

\section{What is H2O?}
\Urlmuskip=0mu plus 1mu\relax %needed to make long URLs break nicely


H2O is fast, scalable, open-source machine learning and deep learning for smarter applications. With H2O, enterprises like PayPal, Nielsen Catalina, Cisco, and others can use all their data without sampling to get accurate predictions faster. Advanced algorithms such as deep learning, boosting, and bagging ensembles are built-in to help application designers create smarter applications through elegant APIs. Some of our initial customers have built powerful domain-specific predictive engines for recommendations, customer churn, propensity to buy, dynamic pricing, and fraud detection for the insurance, healthcare, telecommunications, ad tech, retail, and payment systems industries.

Using in-memory compression, H2O handles billions of data rows in-memory, even with a small cluster. To make it easier for non-engineers to create complete analytic workflows, H2O's platform includes interfaces for R, Python, Scala, Java, JSON, and CoffeeScript/JavaScript, as well as a built-in  web interface, Flow. H2O was built alongside (and on top of) Hadoop and Spark Clusters and typically deploys within minutes.

H2O includes many common machine learning algorithms, such as generalized linear modeling (linear regression, logistic regression, etc.), Na\"{i}ve Bayes, principal components analysis, time series, k-means clustering, and others. H2O also implements best-in-class algorithms at scale, such as distributed random forest, gradient boosting and deep learning. Customers can build thousands of models and compare the results to get the best predictions.

H2O is nurturing a grassroots movement of physicists, mathematicians, and computer scientists to herald the new wave of discovery with data science by collaborating closely with academic researchers and Industrial data scientists. Stanford university giants Stephen Boyd, Trevor Hastie, Rob Tibshirani advise the H2O team on building scalable machine learning algorithms. With hundreds of meetups over the past three years, H2O has become a word-of-mouth phenomenon, growing amongst the data community by a hundred-fold, and is now used by 30,000+ users and is deployed using R, Python, Hadoop, and Spark in 2000+ corporations.

\textbf{Try it out}

\begin{itemize}
\item  Download H2O directly at \mbox{\url{http://h2o.ai/download}}.
\item Install H2O's R package from CRAN at {\url{https://cran.r-project.org/web/packages/h2o/}}. 
\item Install the Python package from PyPI at {\url{https://pypi.python.org/pypi/h2o/}}.

\end{itemize}



\textbf{Join the community}
\begin{itemize}
\item  To learn about our meetups, training sessions, hackathons, and product updates, visit {\url{http://h2o.ai}}. 
\item Visit the open source community forum at {\url{https://groups.google.com/d/forum/h2ostream}}.
\item Join the chat at {\url{https://gitter.im/h2oai/h2o-3}}.
\end{itemize}





\section{Introduction}

This documentation describes the functionality of R in H2O. Further information on H2O's system and algorithms (as well as R user documentation) can be found at the H2O website at {\url{http://docs.h2o.ai}}. This introductory section describes how H2O works with R, followed by a brief overview of generalized linear models (GLM). 

R requires a reference object (transparent to user) to the H2O instance because it uses a REST API to send functions to H2O. Data sets are not transmitted directly through the REST API. Instead, the user sends a command (for example, an HDFS path to the data set) either through the browser or via the REST API to ingest data from disk.

The data set is then assigned a Key in H2O that you can use as a reference in future commands to the web server. After preparing your dataset for modeling by defining the significant data and removing the insignificant data, you can create models to represent the results of the data analysis. One of the most popular models for data analysis is GLM. 

GLM estimates regression models for outcomes following exponential distributions in general. In addition to the Gaussian (i.e. normal) distribution, these include Poisson, binomial, gamma and Tweedie distributions. Each serves a different purpose, and depending on distribution and link function choice, it can be used either for prediction or classification.

H2O supports Spark, YARN, and all versions of Hadoop. Hadoop is a scalable open-source file system that uses clusters to enable distributed storage and processing of datasets. Depending on your data size, you can get started on your desktop or scale using multiple nodes with Hadoop. 

H2O nodes run as JVM invocations on Hadoop nodes. For performance reasons, we recommend you avoid running an H2O node on the same hardware as the Hadoop NameNode if possible.

Since H2O nodes run as mapper tasks in Hadoop, administrators can see them in the normal JobTracker and TaskTracker frameworks. This provides process-level (i.e. JVM instance-level) visibility.

H2O helps R users make the leap from laptop-based processing to large-scale environments. Hadoop helps H2O users scale their data processing capabilities based on their current needs. Using H2O, R, and Hadoop, you can create a complete end-to-end data analysis solution. For more information about H2O on Hadoop, refer to {\url{http://docs.h2o.ai}}.

This document will walk you through the four steps of data analysis with H2O: installing H2O, preparing your data for modeling (data munging), creating a model using state-of-the-art machine learning algorithms, and scoring your models. 

\section{Installation}

To use H2O with R, you can start H2O outside of R and connect to it, or you can launch H2O from R. However, if you launch H2O from R and close the R session, the H2O instance is closed as well. The client object is used to direct R to datasets and models located in H2O.

\subsection{Installing R or R Studio}

To download R:
\begin{enumerate}
\item Go to \url{http://cran.r-project.org/mirrors.html}. 
\item Select your closest local mirror. 
\item Select your operating system (Linux, OS X, or Windows). 
\item Depending on your OS, download the appropriate file, along with any required packages. 
\item When the download is complete, unzip the file and install. \\
\end{enumerate}

To download R Studio: 

\begin{enumerate}
\item Go to \url{http://www.rstudio.com/products/rstudio/}. 
\item Select your deployment type (desktop or server). 
\item Download the file. 
\item When the download is complete, unzip the file and install.
\end{enumerate}


\subsection{Installing H2O in R}

Load the latest CRAN H2O package by running \begin{spverbatim} install.packages("h2o") \end{spverbatim} 

Note: Our push to CRAN will be behind the bleeding edge version and due to resource constraints, may be behind the published version. However, there is a best-effort to keep the versions the same. 

\begin{lstlisting}[style=R]
# The following two commands remove any previously installed H2O packages for R.
if ("package:h2o" %in% search()) { detach("package:h2o", unload=TRUE) }
if ("h2o" %in% rownames(installed.packages())) { remove.packages("h2o") }

# Next, we download packages that H2O depends on.
if (! ("methods" %in% rownames(installed.packages()))) { install.packages("methods") }
if (! ("statmod" %in% rownames(installed.packages()))) { install.packages("statmod") }
if (! ("stats" %in% rownames(installed.packages()))) { install.packages("stats") }
if (! ("graphics" %in% rownames(installed.packages()))) { install.packages("graphics") }
if (! ("RCurl" %in% rownames(installed.packages()))) { install.packages("RCurl") }
if (! ("jsonlite" %in% rownames(installed.packages()))) { install.packages("jsonlite") }
if (! ("tools" %in% rownames(installed.packages()))) { install.packages("tools") }
if (! ("utils" %in% rownames(installed.packages()))) { install.packages("utils") }
\end{lstlisting}


\begin{lstlisting}[style=R]
# Now we download, install and initialize the H2O package for R (the
#below url can be obtained from the H2O download page by clicking on Bleeding Edge, or a given release)
install.packages("h2o", type="source", repos=(c("http://h2o-release.s3.amazonaws.com/h2o/rel-simons/7/R")))
library(h2o)
h2o.init(nthreads = -1)
\end{lstlisting}

\subsection{Making a build from the Source Code}
If you are a developer who wants to make changes to the R package before building and installing it, pull the source code from Git ({\url{https://github.com/h2oai/h2o-3}}) and follow the instructions in at {\url{https://github.com/h2oai/h2o-3/blob/master/README.md}}.

After making the build, navigate to the top-level \texttt{h2o-3} directory using {\texttt{cd $\mathtt{\sim}$/h2o-3}}, then run the following (replacing the asterisks [*] with the version number) and install.
\begin{lstlisting}[style=R]
./gradlew clean 
./gradlew build
$ R CMD INSTALL h2o-r/R/src/contrib/h2o_****.tar.gz
* installing to library �/Users/H2OUser/.Rlibrary�
* installing *source* package �h2o� ...
** R
** demo
** inst
** preparing package for lazy loading
Creating a generic function for ...[output truncated]
** help
*** installing help indices
** building package indices
** testing if installed package can be loaded
* DONE (h2o)
\end{lstlisting}

\section{H2O Initialization}

This section describes how to launch H2O: 
\begin{itemize}
\item from R
\item from the command line
\item on Hadoop
\item on an EC2 cluster
\end{itemize}

\subsection{Launching from R} \label{ssec:LaunchR}

If you do not specify the argument {\texttt{max\_mem\_size}} when you run {\texttt{h2o.init(nthreads = -1)}}, the default heap size of the H2O instance running on 32-bit Java is 1g. H2O checks the Java version and suggests an upgrade if you are running 32-bit Java. On 64-bit Java, the heap size is 1/4 of the total memory available on the machine. The {\texttt{nthreads = -1}} parameter allows H2O to use all CPUs on the host, which is recommended (*mal*the default is two).


For best performance, the allocated memory should be 4x the size of your data, but never more than the total amount of memory on your computer. For larger data sets, we recommend running on a server or service with more memory available for computing.


To launch H2O locally from R, you can run the following in R:
\begin{lstlisting}[style=R]
library(h2o)
# Start H2O on localhost, port 54321, with 4g of memory using all CPUs
h2o.init(ip = 'localhost', port = 54321, nthreads= -1, max_mem_size = '4g')
\end{lstlisting}

R displays the following output: 
\begin{lstlisting}[style=R]
Successfully connected to http://localhost:54321
R is connected to H2O cluster:
   H2O cluster uptime:         11 minutes 35 seconds
   H2O cluster version:        2.7.0.1497
   H2O cluster name:           H2O_started_from_R
   H2O cluster total nodes:    1
   H2O cluster total memory:   3.56 GB
   H2O cluster total cores:    8
   H2O cluster allowed cores:  8
   H2O cluster healthy:        TRUE
\end{lstlisting}

You can also launch H2O locally (with default initiallization arguments) as follows:
\begin{lstlisting}[style=R]
h2o.init()
\end{lstlisting}

H2O may have been launched externally (on a remote server, for example). You may still connect to this H2O from a local R session by specifying the H2O's IP address and port number as follows:
\begin{lstlisting}[style=R]
h2o.init(ip = "192.555.1.123", port = 12345, nthreads = -1)
\end{lstlisting}


\subsection{Launching from the Command Line}

A simple way to launch H2O from the command line is to download the H2O zip file from the H2O download page. Unzip and
launch H2O with the following (replace "h2o-3.3.0.3177.zip" with your respective zip file):
\begin{lstlisting}[style=R]
unzip h2o-3.3.0.3177.zip
cd h2o-3.3.0.3177
java -jar h2o.jar
\end{lstlisting}
See the H2O Documentation for additional JVM and H2O command line options.
After launching the H2O instance, connect to it from R with {\texttt{h2o.init()}} as described above. Again, if H2O was launched remotely
connect to it by specifying the IP address and port number of the H2O instance.

\subsection{Launching on Hadoop}

To launch H2O nodes and form a cluster on the Hadoop cluster, run:
\begin{lstlisting}[style=R]
hadoop jar h2odriver.jar -nodes 1 -mapperXmx 6g -output hdfsOutputDirName
\end{lstlisting}

\begin{itemize}
\item The user will need to launch the Hadoop-specific H2O driver jar (h2odriver.jar) that corresponds to their distribution of Hadoop. Specific driver jar files are available for the following Hadoop versions {\texttt{cdh5.2, cdh5.3, cdh5.4.2, hdp2.1, hdp2.2, mapr3.1.1, mapr4.0.1, mapr5.0}}.
\item The above command launches exactly one 6g node of H2O; however,  we recommend launching the cluster with 4 times the memory of your data file.
\item{\texttt{mapperXmx}} is the mapper size or the amount of memory allocated to each node.
\item{\texttt{nodes}} is the number of nodes requested to form the cluster.
\item{\texttt{output}} is the name of the directory created each time a H2O cloud is created so it is necessary for the name to be unique each time it is launched.
\end{itemize}


\subsection{Launching on an EC2}

{\textbf{Note}}: If you would like to try out H2O on an EC2 cluster, {\url{http://play.h2o.ai}} is the easiest way to get started. H2O Play provides access to a temporary cluster managed by H2O. 

If you would still like to set up your own EC2 cluster, follow the instructions below to build a cluster of EC2 instances by running the following commands on the host that can access the nodes using a public DNS name. 

\begin{enumerate}

\item Edit {\texttt{h2o-cluster-launch-instances.py}}to include your SSH key name and security group name, as well as any other environment-specific variables: 

\begin{lstlisting}[style=Python]
./h2o-cluster-launch-instances.py
./h2o-cluster-distribute-h2o.sh
\end{lstlisting}

 ---OR---

\begin{lstlisting}[style=Python]
./h2o-cluster-launch-instances.py
./h2o-cluster-download-h2o.sh
\end{lstlisting}

{\textbf{Note}}: The second method may be faster than the first, since download pulls from S3. 

\item Distribute the credentials using {\texttt{./h2o-cluster-distribute-aws-credentials.sh.}}

{\textbf{Note}}: If you are running H2O using an IAM role, it is not necessary to distribute the AWS credentials to all the nodes in the cluster. The latest version of H2O can access the temporary access key. 

{\textbf{Caution}}: Distributing the AWS credentials copies the Amazon {\texttt{AWS\_ACCESS\_KEY\_ID}} and {\texttt{AWS\_SECRET\_ACCESS\_KEY}} to the instances to enable S3 and S3N access. Use caution when adding your security keys to the cloud. 

\item Start H2O by launching one H2O node per EC2 instance: {\texttt{./h2o-cluster-start-h2o.sh}}. Wait 60 seconds after entering the command before entering it on the next node. 

\item In your internet browser, substitute any of the public DNS node addresses for {\texttt{IP\_ADDRESS}} in the following example: {\texttt{http://IP\_ADDRESS:54321}}

\item To start H2O: {\texttt{./h2o-cluster-start-h2o.sh}}

\item To stop H2O: {\texttt{./h2o-cluster-stop-h2o.sh}}

\item To shut down the cluster, use your Amazon AWS console to shut down the cluster manually. For more information, refer to the Amazon AWS documentation at {\url{http://docs.aws.amazon.com/ElasticMapReduce/latest/DeveloperGuide/UsingEMR_TerminateJobFlow.html}}. 

\end{enumerate}

\subsection{Checking Cluster Status}


To check the status and health of the H2O cluster, use {\texttt{h2o.clusterInfo()}}.
\begin{lstlisting}[style=R]
> library(h2o)
> h2o.init()
> h2o.clusterInfo()
\end{lstlisting}

An easy-to-read summary of information about the cluster displays. 
\begin{lstlisting}[style=R]
R is connected to H2O cluster:
  H2O cluster uptime:         43 minutes 43 seconds
  H2O cluster version:        2.7.0.1497
  H2O cluster name:           H2O_started_from_R
  H2O cluster total nodes:    1
  H2O cluster total memory:   3.56 GB
  H2O cluster total cores:    8
  H2O cluster allowed cores:  8
  H2O cluster healthy:        TRUE
\end{lstlisting}


\section{Data Preparation in R}

The following section describes some important points to remember about data preparation (also known as data munging) and some of the tools and methods available in H2O, as well as a data training example. 

\subsection{Notes}
\begin{itemize}

\item Although it may seem like you are manipulating the data in R due to the look and feel, once the data has been passed to H2O, all data munging occurs in the H2O instance. The information is passed to R through JSON APIs, so some functions may not have another method. 
\item You are not limited by R's ability to handle data, but by the total amount of memory allocated to the H2O instance. To process large data sets, make sure to allocate enough memory. For more information, refer to \nameref{ssec:LaunchR}. 
\item You can manipulate datasets with thousands of factor levels using H2O in R, so if you ask H2O to display a table in R with information from high cardinality factors, the results may overwhelm R`s capacity. 
\item To manipulate data in R and not in H2O, use {\texttt{as.data.frame()}}, {\texttt{as.h2o()}}, and {\texttt{str()}}. \begin{itemize}
\item {\texttt{as.data.frame()}} converts an H2O data frame into an R data frame. Be aware that if your request exceeds R$'$s capabilities due to the amount of data, the R session will crash. If possible, we recommend only taking subsets of the entire data set (the necessary data columns or rows), and not the whole data set. 
\item {\texttt{as.h2o()}} transfers data from R to the H2O instance. We recommend ensuring that you allocate enough memory to the H2O instance for successful data transfer.
\item {\texttt{str()}} returns the elements of the new object to confirm that the data transferred correctly. It$'$s a good way to verify there were no data loss or conversion issues. %% is this still supported? can't find it in the R package doc...
\end{itemize}
\end{itemize}

\subsection{Tools and Methods}

The following section describes some of the tools and methods available in H2O for data preparation. 
\begin{itemize}
\item {\textbf{Data Profiling}}: Quickly summarize the shape of your dataset to avoid bias or missing information before you start building your model. Missing data, zero values, text, and a visual distribution of the data are visualized automatically upon data ingestion. 
\item {\textbf{Summary Statistics}}: Visualize your data with summary statistics to get the mean, standard deviation, min, max, or quantile (for numeric columns) or cardinality and counts (for enum columns), and a preview of the data set. 
\item {\textbf{Aggregate, Filter, Bin, and Derive Columns}}: Build unique views with Group functions, Filtering, Binning, and Derived Columns. 
\item {\textbf{Slice, Log Transform, and Anonymize}}: Normalize and partition to get your data into the right shape for modeling, and anonymize to remove confidential information. 
\item {\textbf{Variable Creation}}: Highly customizable variable value creation to hone in on the key data characteristics to model. 
\item {\textbf{PCA}}: Principal Component Analysis makes feature selection easy with a simple interface and standard input values to reduce the many dimensions in your dataset into key components. 
\item {\textbf{Training and Validation Sampling Plan}}: Design a random or stratified sampling plan to generate data sets for model training and scoring. 
\end{itemize}

\subsection{Demo: Creating Aggregates from Split Data}

The following section depicts an example of creating aggregates for data training using {\texttt{ddply()}}. Using this method, you can split your dataset and apply a function to the subsets.

To apply a user-specified function to each subset of an H2O dataset and combine the results, use {\texttt{ddply()}}, with the name of the H2O object, the variable name, and the function in the parentheses.

\begin{lstlisting}[style=R]
> library(h2o)
> h2o.init(nthreads = -1)

# Import iris dataset to H2O
> irisPath = system.file("extdata", "iris_wheader.csv", package = "h2o")
> iris.hex = h2o.importFile(path = irisPath, destination_frame = "iris.hex")

# Apply function to groups by class of flower
# uses h2o's ddply, since iris.hex is an H2OFrame object
> res = h2o.ddply(iris.hex, "class", function(df) { sum(df[,1], na.rm = T)/nrow(df) })
> head(res)
\end{lstlisting}

\section{Models}

The following section describes the features and functions of some common models available in H2O.  For more information about running these models in R using H2O, refer to ``Running Models." 

%%mal H2O supports the following models: Deep Learning (DL), Generalized Linear Models (GLM), Gradient Boosted Regression (GBM), Na\"{i}ve Bayes (NB), Random Forest (RF), K-Means, Principal Components Analysis (PCA) and Generalized Low Rank Model (GLRM).
H2O supports the following models:  Deep Learning, Generalized Linear Models (GLM), Gradient Boosted Regression (GBM), Na\"{i}ve Bayes (NB), Random Forest (RF), K-Means, Principal Components Analysis (PCA) and Generalized Low Rank Model (GLRM).

The list is growing quickly, so check back often at \url{www.h2o.ai} to see the latest additions. The following list describes some common model types and features. 

{\textit{Supervised Learning}}

{\textbf{Generalized Linear Models (GLM)}}: A flexible generalization of ordinary linear regression for response variables that have error distribution models other than a normal distribution. GLM unifies various other statistical models, including linear, logistic, Poisson, and more with L1 and L2 regularization.

%%mal {\textbf{Decision trees}}: Used in RF; a decision support tool that uses a tree-like graph or model of decisions and their possible consequences.
{\textbf{Random Forest}}: A method of averaging together multiple decision trees, each created on different random samples of rows and columns. It is one of the most robust algorithms to noisy data, easy to use, non-linear, and provides feedback on the importance of each predictor in the model.

{\textbf{Gradient Boosting (GBM)}}: A method to produce a prediction model in the form of an ensemble of weak prediction models. It builds the model in a stage-wise fashion and is generalized by allowing an arbitrary differentiable loss function. It is one of the most powerful methods available today.

{\textbf{Deep Learning}}: Model high-level abstractions in data by using non-linear transformations in a layer-by-layer method. Deep learning is an example of supervised learning and can make use of unlabeled data that other algorithms cannot.

{\textbf{Na\"{i}ve Bayes}}: A probabilistic classifier that assumes the value of a particular feature is unrelated to the presence or absence of any other feature, given the class variable. It is often used in text categorization.

{\textit{Unsupervised Learning}}

{\textbf{K-Means}}: A method to uncover groups or clusters of data points often used for segmentation. It clusters observations into k certain points with the nearest mean.

{\textbf{Anomaly Detection}}: Identify the outliers in your data by invoking a powerful pattern recognition model, the Deep Learning Auto-Encoder.

{\textit{Modeling Constructs}}

{\textbf{Grid Search}}: The standard way of performing hyper-parameter optimization to make model configuration easier. %%should be supported in latest version, but adding note in case

After creating a model, use it to make predictions. For more information about predictions, refer to ``Predictions." 

%%mal: this feels too lengthy for a modeling example here; perhaps it would be better to have entire work flows at the end?
\subsection{Demo: GLM}

The following demo demonstrates how to import a file, define significant data, view data, create testing and training sets using sampling, define the model, and display the results.

\begin{lstlisting}[style=R]
# Import dataset and display summary
> library(h2o)
> h2o.init()
> airlinesURL = "https://s3.amazonaws.com/h2o-airlines-unpacked/allyears2k.csv"
> airlines.hex = h2o.importFile(path = airlinesURL, destination_frame = "airlines.hex")
> summary(airlines.hex)

# View quantiles and histograms
#high_na_columns = h2o.ignoreColumns(data = airlines.hex)
> quantile(x = airlines.hex$ArrDelay, na.rm = TRUE)
> h2o.hist(airlines.hex$ArrDelay)

# Find number of flights by airport
> originFlights = h2o.ddply(airlines.hex, 'Origin', nrow)
> originFlights.R = as.data.frame(originFlights)

# Find number of cancellations per month
> flightsByMonth = h2o.ddply(airlines.hex,"Month", nrow)
> flightsByMonth.R = as.data.frame(originFlights)

# Find months with the highest cancellation ratio
> fun = function(df) {sum(df[,which(colnames(airlines.hex)=="Cancelled")])}
> cancellationsByMonth = h2o.ddply(airlines.hex,"Month", fun)
> cancellation_rate = cancellationsByMonth$C1/flightsByMonth$C1
> rates_table = cbind(flightsByMonth$Month, cancellation_rate)
> rates_table.R = as.data.frame(rates_table)

# Construct test and train sets using sampling
> airlines.split = h2o.splitFrame(data = airlines.hex,ratios = 0.85)
> airlines.train = airlines.split[[1]]
> airlines.test = airlines.split[[2]]

# Display a summary using table-like functions
> h2o.table(airlines.train$Cancelled)
> h2o.table(airlines.test$Cancelled)

# Set predictor and response variables
> Y = "IsDepDelayed"
> X = c("Origin", "Dest", "DayofMonth", "Year", "UniqueCarrier", "DayOfWeek", "Month", "DepTime", "ArrTime", "Distance")
# Define the data for the model and display the results
> airlines.glm <- h2o.glm(training_frame=airlines.train, x=X, y=Y, family = "binomial", alpha = 0.5)
# View model information: training statistics, performance, important variables
> summary(airlines.glm)

# Predict using GLM model
> pred = h2o.predict(object = airlines.glm, newdata = airlines.test)
# Look at summary of predictions: probability of TRUE class (p1)
> summary(pred$p1)

\end{lstlisting}




\section{Data Manipulation in R}

The following section describes some common R commands. For a complete command list, including parameters, refer to {\url{http://h2o-release.s3.amazonaws.com/h2o/latest_stable_Rdoc.html}}.
For additional help within R's Help tab, precede the command with a question mark (for example, {\texttt{?h2o}}) for suggested commands containing the search terms. For more information on a command, precede the command with two question marks ({\texttt{??h2o}}). 

\subsection{Importing Files}

The H2O package consolidates all of the various supported import functions using {\texttt{h2o.importFile()}}. Although {\texttt{h2o.importFolder}} and {\texttt{h2o.importHDFS}} will still work, these functions are deprecated and should be updated to {\texttt{h2o.importFile()}}. There are a few ways to import files: 

\begin{lstlisting}[style=R]

# To import small iris data file from H2O's package:
> irisPath = system.file("extdata", "iris.csv", package="h2o")
> iris.hex = h2o.importFile(path = irisPath, destination_frame = "iris.hex")
|=================================================| 100%

# To import an entire folder of files as one data object:
> pathToFolder = "/Users/Amy/data/airlines/"
> airlines.hex = h2o.importFile(path = pathToFolder, destination_frame = "airlines.hex")
|=================================================| 100%

# To import from HDFS, connect to H2O in R using the IP and port of an H2O instance running on your Hadoop cluster:
> h2o.init(ip= <IPAddress>, port =54321, nthreads = -1)
> pathToData = "hdfs://mr-0xd6.h2oai.loc/datasets/airlines_all.csv"
> airlines.hex = h2o.importFile(path = pathToData, destination_frame = "airlines.hex")
|=================================================| 100%
\end{lstlisting}


\subsection{Uploading Files}

To upload a file in a directory local to your H2O instance, we recommend {\texttt{h2o.importFile()}}. The alternative is to use {\texttt{h2o.uploadFile()}} which can also upload data local to your H2O instance in addition to uploading data local to your R session. In the parentheses, specify the H2O reference object in R and the complete URL or normalized file path for the file.
\begin{lstlisting}[style=R]
irisPath = system.file("extdata", "iris.csv", package="h2o")
iris.hex = h2o.uploadFile(path = irisPath, destination_frame = "iris.hex")
|====================================================| 100% 
\end{lstlisting}


\subsection{Finding Factors}

To determine if any column in a data set is a factor (contains categorical data), use {\texttt{h2o.anyFactor()}} with the name of the R reference object in the parentheses.
\begin{lstlisting}[style=R]
> irisPath = system.file("extdata", "iris_wheader.csv", package="h2o")
> iris.hex = h2o.importFile(path = irisPath)
|===================================================| 100%
> h2o.anyFactor(iris.hex)
[1] TRUE
\end{lstlisting}

\subsection{Converting to Factors}
To convert an integer into a non-ordered factor (also called an enum or categorical), use {\texttt{as.factor()}} with the name of the R reference object in parentheses followed by the number of the column to convert in brackets.
\begin{lstlisting}[style=R]
# Import prostate data
> prosPath <- system.file("extdata", "prostate.csv", package="h2o")
> prostate.hex <- h2o.importFile(path = prosPath)
|===================================================| 100%
# Converts column 4 (RACE) to an enum
> is.factor(prostate.hex[,4])
[1] FALSE
> prostate.hex[,4] < -as.factor(prostate.hex[,4])
> is.factor(prostate.hex[,4])
[1] TRUE
# Summary will return a count of the factors
> summary(prostate.hex[,4])
 RACE   
 1 :341 
 2 : 36 
 0 :  3 
\end{lstlisting}
\subsection{Converting Data Frames}

To convert an H2O parsed data object into an R data frame that can be manipulated using R commands, use {\texttt{as.data.frame()}} with the name of the R reference object in the parentheses.

{\textbf{Caution}}: While this can be very useful, be careful using this command when converting H2O parsed data objects. H2O can easily handle data sets that are often too large to be handled equivalently well in R. 
\begin{lstlisting}[style=R]
# Creates object that defines path
prosPath <- system.file("extdata", "prostate.csv", package="h2o")
# Imports data set
prostate.hex = h2o.importFile(path = prosPath)
|===================================================| 100%
# Converts current data frame (prostate data set) to an R data frame
prostate.R <- as.data.frame(prostate.hex)
# Displays a summary of data frame where the summary was executed in R
summary(prostate.R) 
       ID            CAPSULE            AGE             RACE
Min.   :  1.00   Min.   :0.0000   Min.   :43.00   Min.   :0.000
1st Qu.: 95.75   1st Qu.:0.0000   1st Qu.:62.00   1st Qu.:1.000
       .... 
\end{lstlisting}


\subsection{Transferring Data Frames}
To transfer a data frame from the R environment to the H2O instance, use  {\texttt{as.h2o()}}. In the parentheses, specify the object in the R environment to be converted to an H2O object. Optionally, you can include the name of the destination frame in H2O. Precede the destination frame name with {\texttt{destination_frame=}} and enclose the name in quotes as in the following example.

\begin{lstlisting}[style=R]
# Import the iris data into H2O
> data(iris)
> iris
    Sepal.Length Sepal.Width Petal.Length Petal.Width    Species
1            5.1         3.5          1.4         0.2     setosa
2            4.9         3.0          1.4         0.2     setosa
3            4.7         3.2          1.3         0.2     setosa
4            4.6         3.1          1.5         0.2     setosa
5            5.0         3.6          1.4         0.2     setosa
6            5.4         3.9          1.7         0.4     setosa

# Converts R object "iris" into H2O object�"iris.hex"
> iris.hex = as.h2o(iris, destination_frame= "iris.hex")
|=============================================================| 100%
> head(iris.hex)
  Sepal.Length Sepal.Width Petal.Length Petal.Width Species
1          5.1         3.5          1.4         0.2  setosa
2          4.9         3.0          1.4         0.2  setosa
3          4.7         3.2          1.3         0.2  setosa
4          4.6         3.1          1.5         0.2  setosa
5          5.0         3.6          1.4         0.2  setosa
6          5.4         3.9          1.7         0.4  setosa
\end{lstlisting}


\subsection{Renaming Data Frames}

To rename a dataframe on the server running H2O for a data set manipulated in R, use {\texttt{h2o.assign()}}. For instance, in the following example, the prostate data set was uploaded to the H2O instance and the data was manipulated to remove outliers. {\texttt{h2o.assign()}} saves the new data set on the H2O server so that it can be analyzed using H2O without overwriting the original data set.  

\begin{lstlisting}[style=R]
prosPath <- system.file("extdata", "prostate.csv", package="h2o")
prostate.hex<-h2o.importFile(path = prosPath)
|===================================================| 100%
## Assign a new name to prostate dataset in the KV store
prostate.hex@frame_id
[1] "prostate.hex"
prostate.hex <- h2o.assign(data = prostate.hex, key = "prostate.hex")
prostate.hex@frame_id
[1] "prostate.hex"

\end{lstlisting}

\subsection{Getting Column Names}

To obtain a list of the column names in the data set, use {\texttt{colnames()}} or {\texttt{names()}} with the name of the R reference object in the parentheses.

\begin{lstlisting}[style=R]
##Displays the titles of the columns
> colnames(iris.hex)
[1] "Sepal.Length" "Sepal.Width"  "Petal.Length" "Petal.Width"  "Species"     
> names(iris.hex)
[1] "Sepal.Length" "Sepal.Width"  "Petal.Length" "Petal.Width"  "Species"     
\end{lstlisting}


\subsection{Getting Minimum and Maximum Values}

To obtain the maximum values for the real-valued columns in a data set, use {\texttt{max()}} with the name of the R reference object in the parentheses.

To obtain the minimum values for the real-valued columns in a data set, use {\texttt{min()}} with the name of the R reference object in the parentheses.
\begin{spverbatim}
> min(prostate.hex$AGE)
[1] 43
> max(prostate.hex$AGE)
[1] 79
\end{spverbatim}


\subsection{Getting Quantiles}

To request quantiles for an H2O parsed data set, use {\texttt{quantile()}} with the name of the R reference object in the parentheses.
To request a quantile for a single numerical column, use {\texttt{quantile(ReferenceObject\$ColumnName)}},  where {\texttt{ReferenceObject}} represents the R reference object name and {\texttt{ColumnName}} represents the name of the specified column. 
When you request for a full parsed data set consisting of a single column, {\texttt{quantile()}} displays a matrix with quantile information for the data set. 

\begin{lstlisting}[style=R]
> prosPath <- system.file("extdata", "prostate.csv", package="h2o")
> prostate.hex <- h2o.importFile(path = prosPath)
# Returns the percentiles at 0, 10, 20, ..., 100%
> prostate.qs <- quantile(prostate.hex$PSA, probs = (1:10)/10)
> prostate.qs
   10%    20%    30%    40%    50%    60%    70%    80%    90%   100% 
   2.60    4.48    5.77    7.40    8.75   11.00  13.70   20.16   33.21  139.70 
# Take the outliers or the bottom and top 10% of data
> PSA.outliers <- prostate.hex[prostate.hex$PSA <= prostate.qs["10%"] | prostate.hex$PSA >=   prostate.qs["90%"],]
# Check that the number of rows return is about 20% of the original data
> nrow(prostate.hex)
[1] 380
> nrow(PSA.outliers)
[1] 78
> nrow(PSA.outliers)/nrow(prostate.hex)
[1] 0.2052632
\end{lstlisting}


\subsection{Summarizing Data}
To generate a summary (similar to the one in R) for each of the columns in the data set, use {\texttt{summary()}} with the name of the R reference object in the parentheses.
For continuous real functions, this produces a summary that includes information on quartiles, min, max, and mean. 
For factors, this produces information about counts of elements within each factor level. 

\begin{spverbatim}
> summary(prostate.hex)
 ID               CAPSULE          AGE             RACE            DPROS          
 Min.   :  1.00   Min.   :0.0000   Min.   :43.00   Min.   :0.000   Min.   :1.000  
 1st Qu.: 95.75   1st Qu.:0.0000   1st Qu.:62.00   1st Qu.:1.000   1st Qu.:1.000  
 Median :190.50   Median :0.0000   Median :67.00   Median :1.000   Median :2.000  
 Mean   :190.50   Mean   :0.4026   Mean   :66.04   Mean   :1.087   Mean   :2.271  
 3rd Qu.:285.25   3rd Qu.:1.0000   3rd Qu.:71.00   3rd Qu.:1.000   3rd Qu.:3.000  
 Max.   :380.00   Max.   :1.0000   Max.   :79.00   Max.   :2.000   Max.   :4.000  
 DCAPS           PSA               VOL             GLEASON        
 Min.   :1.000   Min.   :  0.300   Min.   : 0.00   Min.   :0.000  
 1st Qu.:1.000   1st Qu.:  4.900   1st Qu.: 0.00   1st Qu.:6.000  
 Median :1.000   Median :  8.664   Median :14.20   Median :6.000  
 Mean   :1.108   Mean   : 15.409   Mean   :15.81   Mean   :6.384  
 3rd Qu.:1.000   3rd Qu.: 17.063   3rd Qu.:26.40   3rd Qu.:7.000  
 Max.   :2.000   Max.   :139.700   Max.   :97.60   Max.   :9.000
 \end{spverbatim}


\subsection{Summarizing Data in a Table}

To summarize the data, use {\texttt{h2o.table()}}. Because H2O can handle larger data sets, it is possible to generate tables that are larger than R`s capacity. 

To summarize multiple columns, use {\texttt{head(h2o.table (ObjectName[, c(ColumnNumber,ColumnNumber)]))}} where {\texttt{ObjectName}} is the name of the object in R and {\texttt{ColumnNumber}} is the number of the column. 

\begin{lstlisting}[style=R]
# Counts of the ages of all patients
> head(as.data.frame(h2o.table(prostate.hex[,"AGE"])))
   AGE Count
1   43     1
2   47     1
3   50     2
4   51     3
5   52     2
6   53     4

# Two-way table of ages (rows) and race (cols) of all patients
# Example: For the first row there is one count of a 43 year old that's labeled as RACE = 0
> h2o.table(prostate.hex[,c("AGE","RACE")])
H2OFrame with 53 rows and 3 columns

First 10 rows:
   AGE RACE count
1   53    1     3
2   61    1    12
3   70    0     1
4   75    1    11
5   74    1    13
6   76    2     1
7   53    2     1
8   52    1     2
9   61    2     1
10  60    1     9
\end{lstlisting}

\subsection{Generating Random Uniformly Distributed Numbers}

To append a column of random numbers to an H2O data frame for facilitating creation of testing/training data splits for analysis and validation in H2O, use {\texttt{h2o.runif()}} with the name of the R reference object in the parentheses. This method is best for customized frame splitting; otherwise, use {\texttt{h2o.splitFrame()}}. However, {\texttt{h2o.runif()}} is not as fast or stable as {\texttt{h2o.splitFrame()}}. 

\begin{lstlisting}[style=R]
> prosPath <- system.file("extdata", "prostate.csv", package="h2o")
> prostate.hex <- h2o.importFile(path = prosPath)

## Creates object for uniform distribution on prostate data set
> s <- h2o.runif(prostate.hex)
> summary (s)  ## Summarize the results of h2o.runif
 rnd               
 Min.   :0.000863  
 1st Qu.:0.239763  
 Median :0.507936  
 Mean   :0.506718  
 3rd Qu.:0.765194  
 Max.   :0.993178  
## Create training set with threshold of 0.8
> prostate.train <- prostate.hex[s <= 0.8,]
##Assign name to training set
> prostate.train <- h2o.assign(prostate.train, "prostate.train")
## Create test set with threshold to filter values greater than 0.8
> prostate.test <- prostate.hex[s > 0.8,]
## Assign name to test set
> prostate.test <- h2o.assign(prostate.test, "prostate.test")
## Combine results of test & training sets, then display result
> nrow(prostate.train) + nrow(prostate.test)
[1] 380
> nrow(prostate.hex) ## Matches the full set
[1] 380
\end{lstlisting}
 

\subsection{Splitting Frames}

To generate two subsets (according to specified ratios) from an existing H2O data set for testing/training, use {\texttt{h2o.splitFrame()}}.  h2o.splitFrame() returns contiguous sections of the data (it does not randomly sample).

\begin{spverbatim}
# Splits data in prostate data frame with a ratio of 0.75
> prostate.split <- h2o.splitFrame(data = prostate.hex , ratios = 0.75)
# Creates training set from 1st data set in split
> prostate.train <- prostate.split[[1]]
# Creates testing set from 2st data set in split
> prostate.test <- prostate.split[[2]]
\end{spverbatim}


\subsection{Getting Frames}

To create a reference object to the data frame in H2O, use {\texttt{h2o.getFrame()}}. This is helpful for  users that alternate between the web UI and the R API or multiple users accessing the same H2O instance. The following example assumes prostate.hex is in the key-value (KV) store.

\begin{spverbatim}
> prostate.hex <- h2o.getFrame(frame_id = "prostate.hex")
\end{spverbatim}


\subsection{Getting Models}
To create a reference object for the model in H2O, use {\texttt{h2o.getModel()}}. This is helpful for  users that alternate between the web UI and the R API or multiple users accessing the same H2O instance. The following example assumes GBM_8e4591a9b413407b983d73fbd9eb44cf is in the key-value (KV) store.
\begin{spverbatim}
> gbm.model <- h2o.getModel(model_id = "GBM_8e4591a9b413407b983d73fbd9eb44cf")
\end{spverbatim}

\subsection{Listing H2O Objects}

To generate a list of all H2O objects generated during a session, along with each object’s size in bytes, use {\texttt{h2o.ls()}}.

\begin{spverbatim}
> h2o.ls()
                                                   Key       Bytesize
      1               GBM_8e4591a9b413407b983d73fbd9eb44cf    40617
      2               GBM_a3ae2edf5dfadbd9ba5dc2e9560c405d     1516
\end{spverbatim}


\subsection{Removing H2O Objects}

To remove an H2O object on the server associated with an object in the R environment, use {\texttt{h2o.rm()}}. For optimal performance, we recommend removing the object from the R environment as well using {\texttt{remove()}}, with the name of the object in the parentheses. If you do not specify an R environment, then the current environment is used. 
\begin{spverbatim}
> h2o.rm(ids = c("prostate.train","prostate.test"))
> h2o.ls()
\end{spverbatim}


\subsection{Adding Functions}

User-defined functions no longer need to be added explicitly to the H2O instance. An R function can be defined and executed against an H2OFrame. 
\begin{spverbatim}
# Create an R functional expression
> simpleFun <- function(x) { 2*x + 5 }
# Evaluate the expression across prostate's AGE column
> calculated <- simpleFun(prostate.hex[,"AGE"])
> h2o.cbind(prostate.hex[,"AGE"], calculated)

H2OFrame with 380 rows and 2 columns

First 10 rows:
   AGE AGE0
1   65  135
2   72  149
3   70  145
4   76  157
5   69  143
6   71  147
7   68  141
8   61  127
9   69  143
10  68  141
\end{spverbatim}


\section{Running Models}

\subsection{Gradient Boosted Models (GBM)}
To generate gradient boosted models for developing forward-learning ensembles, use {\texttt{h2o.gbm()}}.  In the parentheses, define x (the predictor variable vector), y (the integer or categorical response variable), the distribution type (multinomial is the default, gaussian is used for regression), and the name of the H2OParsedData object. 

For more information, use {\texttt{help(h2o.gbm)}}.
\begin{spverbatim}
> library(h2o)
> h2o.init(nthreads = -1)
> data(iris)
> iris.hex <- as.h2o(iris,destination_frame = "iris.hex")
> iris.gbm <- h2o.gbm(y = 1, x = 2:5, training_frame = iris.hex, ntrees = 10,
    max_depth = 3,min_rows = 2, learn_rate = 0.2, distribution= "gaussian")

# To obtain the Mean-squared Error by tree from the model object:
> iris.gbm@model$scoring_history

Scoring History:
             timestamp   duration number_of_trees training_MSE training_deviance
1  2015-09-11 09:50:16  0.005 sec               1      0.47256           0.47256
2  2015-09-11 09:50:16  0.008 sec               2      0.33494           0.33494
3  2015-09-11 09:50:16  0.011 sec               3      0.24291           0.24291
4  2015-09-11 09:50:16  0.014 sec               4      0.18414           0.18414
5  2015-09-11 09:50:16  0.017 sec               5      0.14363           0.14363
6  2015-09-11 09:50:16  0.020 sec               6      0.11677           0.11677
7  2015-09-11 09:50:16  0.023 sec               7      0.09916           0.09916
8  2015-09-11 09:50:16  0.026 sec               8      0.08649           0.08649
9  2015-09-11 09:50:16  0.029 sec               9      0.07761           0.07761
10 2015-09-11 09:50:16  0.032 sec              10      0.07071           0.07071

\end{spverbatim}

To generate a classification model that uses labels, use {\texttt{distribution= "multinomial"}}: 
\begin{spverbatim}
> iris.gbm2 <- h2o.gbm(y = 5, x = 1:4, training_frame = iris.hex, ntrees = 15
                     , max_depth = 5, min_rows = 2, learn_rate = 0.01, distribution= "multinomial")
   
> iris.gbm2@model$training_metrics

H2OMultinomialMetrics: gbm
** Reported on training data. **

Training Set Metrics: 
=====================

Extract training frame with `h2o.getFrame("iris.hex")`
MSE: (Extract with `h2o.mse`) 0.3293958
R^2: (Extract with `h2o.r2`) 0.5059063
Logloss: (Extract with `h2o.logloss`) 0.8533637
Confusion Matrix: Extract with `h2o.confusionMatrix(<model>,train=TRUE)`)
=========================================================================
           setosa versicolor virginica      Error    Rate
setosa         50          0         0 0.00000000  0 / 50
versicolor      0         49         1 0.02000000  1 / 50
virginica       0          1        49 0.02000000  1 / 50
Totals         50         50        50 0.01333333 2 / 150

Hit Ratio Table: Extract with `h2o.hit_ratio_table(<model>,train=TRUE)`
=======================================================================
Top-3 Hit Ratios:
  k hit_ratio
1 1  0.986667
2 2  1.000000
3 3  1.000000

\end{spverbatim}

\subsection{Generalized Linear Models (GLM)}


Generalized linear models (GLM) are some of the most commonly-used models for many types of data analysis use cases. While some data analysis can be done using general linear models, if the variables are more complex, general linear models may not be as accurate. For example, if the dependent variable has a non-continuous distribution or if the effect of the predictors is not linear, generalized linear models will produce more accurate results than general linear models.  

Generalized Linear Models (GLM) estimate regression models for outcomes following exponential distributions in general. In addition to the Gaussian (i.e. normal) distribution, these include Poisson, binomial, gamma and Tweedie distributions. Each serves a different purpose, and depending on distribution and link function choice, it can be used either for prediction or classification.

H2O's GLM algorithm fits the generalized linear model with elastic net penalties. The model fitting computation is distributed, extremely fast,and scales extremely well for models with a limited number (\~\ low thousands) of predictors with non-zero coefficients. The algorithm can compute models for a single value of a penalty argument or the full regularization path, similar to glmnet. It can compute gaussian (linear), logistic, poisson, and gamma regression models.


To generate a generalized linear model for developing linear models for exponential distributions, use {\texttt{h2o.glm()}}. You can apply regularization to the model by adjusting the lambda and alpha parameters. 
For more information, use {\texttt{help(h2o.glm)}}.
\begin{spverbatim}
> prostate.hex <- h2o.importFile(path = "https://raw.github.com/h2oai/h2o/master/smalldata/logreg/prostate.csv"
                               , destination_frame = "prostate.hex")

> prostate.glm<-h2o.glm(y = "CAPSULE", x = c("AGE","RACE","PSA","DCAPS"), training_frame = prostate.hex,
              family = "binomial", nfolds = 10, alpha = 0.5)
> prostate.glm@model$cross_validation_metrics

H2OBinomialMetrics: glm
** Reported on cross-validation data. **
Description: 10-fold cross-validation on training data

MSE:  0.2093902
R^2:  0.1294247
LogLoss:  0.6095525
AUC:  0.6909965
Gini:  0.381993
Null Deviance:  513.8229
Residual Deviance:  463.2599
AIC:  473.2599

Confusion Matrix for F1-optimal threshold:
         0   1    Error      Rate
0      122 105 0.462555  =105/227
1       41 112 0.267974   =41/153
Totals 163 217 0.384211  =146/380

Maximum Metrics:
                      metric threshold    value idx
1                     max f1  0.312978 0.605405 216
2                     max f2  0.138305 0.772727 377
3               max f0point5  0.400689 0.628141 110
4               max accuracy  0.400689 0.700000 110
5              max precision  0.998848 1.000000   0
6           max absolute_MCC  0.400689 0.357638 110
7 max min_per_class_accuracy  0.330976 0.621145 181

\end{spverbatim}

\subsection{K-Means}

To generate a K-Means model for data characterization, use {\texttt{h2o.kmeans()}}. This algorithm does not rely on a dependent variable. For more information, use {\texttt{help(h2o.kmeans)}}.
\begin{spverbatim}
> h2o.kmeans(training_frame = iris.hex, k = 3, x = 1:4)

Model Details:
==============

H2OClusteringModel: kmeans
Model ID:  K-means_model_R_1441989204383_30 
Model Summary:
  number_of_rows number_of_clusters number_of_categorical_columns number_of_iterations within_cluster_sum_of_squares
1            150                  3                             0                    8                     139.09920
  total_sum_of_squares between_cluster_sum_of_squares
1            596.00000                      456.90080


H2OClusteringMetrics: kmeans
** Reported on training data. **


Total Within SS:  139.0992
Between SS:  456.9008
Total SS:  596 
Centroid Statistics:
  centroid     size within_cluster_sum_of_squares
1        1 44.00000                      43.34674
2        2 50.00000                      47.35062
3        3 56.00000                      48.40184
\end{spverbatim}

\subsection{Principal Components Analysis (PCA)}

To map a set of variables onto a subspace using linear transformations, use {\texttt{h2o.prcomp()}}. This is the first step in Principal Components Regression. For more information, use {\texttt{help(h2o.prcomp)}}.
\begin{spverbatim}
> ausPath = system.file("extdata", "australia.csv", package="h2o")
> australia.hex = h2o.importFile(path = ausPath)

> australia.pca <- h2o.prcomp(training_frame = australia.hex, transform = "STANDARDIZE",k = 3)
> australia.pca
Model Details:
==============

H2ODimReductionModel: pca
Model Key:  PCA_model_R_1441989204383_36 
Importance of components:
                            pc1      pc2      pc3
Standard deviation     1.750703 1.512142 1.031181
Proportion of Variance 0.383120 0.285822 0.132917
Cumulative Proportion  0.383120 0.668942 0.801859
\end{spverbatim}

\subsection{Predictions}

The following section describes some of the prediction methods available in H2O. 

{\textbf{Predict}}: Generate outcomes of a data set with any model. Predict with GLM, GBM, Decision Trees or Deep Learning models.\bigskip

{\textbf{Confusion Matrix}}: Visualize the performance of an algorithm in a table to understand how a model performs.\bigskip

{\textbf{Area Under Curve (AUC)}}: A graphical plot to visualize the performance of a model by its sensitivity, true positives and false positives to select the best model.\bigskip

{\textbf{Hit Ratio}}: A classification matrix to visualize the ratio of the number of correctly classified and incorrectly classified cases.\bigskip

{\textbf{PCA Score}}: Determine how well your feature selection fits a particular model.\bigskip

{\textbf{Multi-Model Scoring}}: Compare and contrast multiple models on a data set to find the best performer to deploy into production. 

To apply an H2O model to a holdout set for predictions based on model results, use {\texttt{h2o.predict()}}.  
In the following example, H2O generates a model and then displays the predictions for that model.
For classification, the predict column is the model's discrete prediction, based on maximum F1 by default; the individual class probabilities are the remaining columns in the data frame. 
It is common to utilize the p1 column for binary classification, if a raw probability is desired.
\begin{spverbatim}
> prostate.fit = h2o.predict(object = prostate.glm, newdata = prostate.hex)
> prostate.fit

H2OFrame with 380 rows and 3 columns

First 10 rows:
   predict         p0        p1
1        0 0.74476265 0.2552373
2        1 0.39763451 0.6023655
3        1 0.41268532 0.5873147
4        1 0.37270563 0.6272944
5        1 0.64649990 0.3535001
6        1 0.43367145 0.5663285
7        1 0.26542251 0.7345775
8        1 0.06143281 0.9385672
9        0 0.73057373 0.2694263
10       1 0.46709293 0.5329071
\end{spverbatim}


\section{Support}
There are multiple ways to request support for H2O: \bigskip

{\textbf{Email}}: {\url{h2ostream@googlegroups.com}}

{\textbf{H2OStream on Google Groups}}: {\url{https://groups.google.com/d/forum/h2ostream}} 

{\textbf{JIRA}}: {\url{http://jira.0xdata.com/}}

{\textbf{Meetup information}}: {\url{http://h2o.ai/events}}

\section{References}\bigskip

{\bfseries{R Package}}:  \url{http://h2o-release.s3.amazonaws.com/h2o/latest_stable_Rdoc.html}

{\bfseries{R Ensemble documentation}}: \url{http://www.stat.berkeley.edu/~ledell/R/h2oEnsemble.pdf} 

{\bfseries{Slide deck}}: \url{http://h2o.ai/blog/2013/08/big-data-science-in-h2o-with-r/} 

{\bfseries{R project website}}: \url{http://www.r-project.org} 

\newpage
\section{Appendix: Commands} \label{Appendix} 

The following section lists some common commands by function that are available in R and a brief description of each command. 
\subsection {Data Set Operations}
{\emph{Data Import/Export}}\par
{\texttt{h2o.downloadCSV}}: Download a H2O dataset to a CSV file on local disk.\newline
{\texttt{h2o.exportFile}}: Export H2O Data Frame to a file.\newline
{\texttt{h2o.importFile}}: Import a file from the local path and parse it.\newline
{\texttt{h2o.parseRaw}}: Parse a raw data file. \newline
{\texttt{h2o.uploadFile}}: Upload a file from the local drive and parse it.\newline

{\emph{Native R to H2O Coercion}}\par
 {\texttt{as.h2o}}: Convert an R object to an H2O object.\newline


\emph{H2O to Native R Coercion}\par
{\texttt{as.data.frame}}: Check if an object is a data frame, or coerce it if possible.\newline

\emph{Data Generation}\par
{\texttt{h2o.createFrame}}: Create an H2O data frame, with optional randomization.\newline
{\texttt{h2o.runif}}: Produce a vector of random uniform numbers.\newline
{\texttt{h2o.interaction}}: Create interaction terms between categorical features of an H2O Frame.\newline
%%\newpage
\emph{Data Sampling / Splitting}\par
{\texttt{h2o.splitFrame}: Split an existing H2O data set according to user-specified ratios.\newline

\emph{Missing Data Handling}\par
{\texttt{h2o.impute}}: Impute a column of data using the mean, median, or mode.\newline
{\texttt{h2o.insertMissingValues}}: Replaces a user-specified fraction of entries in a H2O dataset with missing values.\newline

\subsection{General Data Operations}

\emph{Subscripting example to pull pieces from data object.} 
\begin{spverbatim} 
  x[j]  ## note: chooses column J, not row J
  x[i, j]
  x[[i]]
  x$name 
  x[i] <- value
  x[i, j, ...] <- value
  x[[i]] <- value
  x$i <- value
\end{spverbatim}

\emph{Subsetting}\newline
{\texttt{head, tail}}: Return the First or Last Part of an Object\newline

\emph{Concatenation}\par
 {\texttt{c}}: Combine Values into a Vector or List\newline%%still supported? couldn't find in R doc
 {\texttt{h2o.cbind}}: Take a sequence of H2O datasets and combine them by column.\par
{\emph{Data Attributes}}\par
{\texttt{colnames}}: Return column names for a parsed H2O data object. \newline
{\texttt{colnames$<$-}}: Retrieve or set the row or column names of a matrix-like object.\newline
{\texttt{names}}: Get the name of an object. \newline
{\texttt{names$<$-}}: Set the name of an object. \newline
{\texttt{dim}}: Retrieve the dimension of an object. \newline
{\texttt{length}}: Get the length of vectors (including lists) and factors. \newline
{\texttt{nrow}}: Return a count of the number of rows in an H2OParsedData object. \newline
{\texttt{ncol}}: Return a count of the number of columns in an H2OParsedData object.\newline
{\texttt{h2o.anyFactor}}: Check if an H2O parsed data object has any categorical data columns. \newline
{\texttt{is.factor}}: Check if a given column contains categorical data.\newline

{\emph{Data Type Coercion}}\par
{\texttt{as.factor}}: Convert a column from numeric to factor.\newline
{\texttt{as.Date}}: Converts a column from factor to date.\newline

\subsection{Methods from Group Generics}

{\emph{Math (H2O)}}\par
{\texttt{abs}}: Compute the absolute value of x. \newline
{\texttt{sign}}: Return a vector with the signs of the corresponding elements of x (the sign of a real number is 1, 0, or -1 if the number is positive, zero, or negative, respectively). \newline
{\texttt{sqrt}}: Computes the principal square root of x, $\sqrt{x}$.\newline
{\texttt{ceiling}}: Take a single numeric argument x and return a numeric vector containing the smallest integers not less than the corresponding elements of x. \newline
{\texttt{floor}}: Take a single numeric argument x and return a numeric vector containing the largest integers not greater than the corresponding elements of x. \newline
{\texttt{trunc}}: Take a single numeric argument x and return a numeric vector containing the integers formed by truncating the values in x toward 0. \newline
{\texttt{log}}: Compute logarithms (by default, natural logarithms). \newline
{\texttt{exp}}: Compute the exponential function.\newline

{\emph{Math (generic)}}\par

{\texttt{cummax}}: Display a vector of the cumulative maxima of the elements of the argument.\newline
{\texttt{cummin}}: Display a vector of the cumulative minima of the elements of the argument.\newline
{\texttt{cumprod}}: Display a vector of the cumulative products of the elements of the argument. \newline
{\texttt{cumsum}}: Display a vector of the cumulative sums of the elements of the argument.  \newline
{\texttt{log10}}: Compute common (i.e., base 10) logarithms \newline
{\texttt{log2}}: Compute binary (i.e., base 2) logarithms. \newline
{\texttt{log1p}}: Compute log(1+x) accurately also for $|x|${}\textless{}\textless{} 1.\newline
{\texttt{acos}}: Compute the trigonometric arc-cosine. \newline
{\texttt{acosh}}: Compute the hyperbolic arc-cosine. \newline
{\texttt{asin}}: Compute the trigonometric arc-sine.\newline
{\texttt{asinh}}: Compute the hyperbolic arc-sine.\newline
{\texttt{atan}}: Compute the trigonometric arc-tangent. \newline
{\texttt{atanh}}: Compute the hyperbolic arc-tangent.\newline
{\texttt{expm1}}: Compute exp(x) - 1 accurately also for $|x|$\textless{}\textless{} 1.\newline
{\texttt{cos}}: Compute the trigonometric cosine.\newline
{\texttt{cosh}}: Compute the hyperbolic cosine. \newline
{\texttt{cospi}}: Compute the trigonometric two-argument arc-cosine. \newline
{\texttt{sin}}: Compute the trigonometric sine. \newline
{\texttt{sinh}}: Compute the hyperbolic sine. \newline
{\texttt{sinpi}}: Compute the trigonometric two-argument arc-sine. \newline
{\texttt{tan}}: Compute the trigonometric tangent.\newline
{\texttt{tanh}}: Compute the hyperbolic tangent.\newline
{\texttt{tanpi}}: Compute the trigonometric two-argument arc-tangent. \newline
{\texttt{gamma}}: Display the gamma function $\gamma{x}$ \newline
{\texttt{lgamma}}: Display the natural logarithm of the absolute value of the gamma function. \newline
{\texttt{digamma}}: Display the first derivative of the logarithm of the gamma function. \newline
{\texttt{trigamma}}: Display the second derivative of the logarithm of the gamma function.\newline

\emph{Math2 (H2O)}\par
{\texttt{round}}: Round the values to the specified number of decimal places (default 0). \newline
{\texttt{signif}}: Round the values to the specified number of significant digits.\newline

\emph{Summary (H2O)}\par
{\texttt{max}}: Display the maximum of all the input arguments. \newline
{\texttt{min}}: Display the minimum of all the input arguments. \newline
{\texttt{range}}: Display a vector containing the minimum and maximum of all the given arguments. \newline
{\texttt{sum}}: Calculate the sum of all the values present in its arguments.\newline

\emph{Summary (generic)}\par
{\texttt{prod}}: Display the product of all values present in its arguments. \newline
{\texttt{any}}: Given a set of logical vectors, determine if at least one of the values is true. \newline
{\texttt{all}}: Given a set of logical vectors, determine if all of the values are true.\newline

\subsection{Other Aggregations}

\emph{Non-Group Generic Summaries}\par
{\texttt{mean}}: Generic function for the (trimmed) arithmetic mean. \newline
{\texttt{sd}}: Calculate the standard deviation of a column of continuous real valued data. \newline
{\texttt{var}}: Compute the variance of x.\newline
{\texttt{summary}}: Produce result summaries of the results of various model fitting functions. \newline
{\texttt{quantile}}: Obtain and display quantiles for H2O parsed data.\newline

\emph{Row / Column Aggregation}\par
{\texttt{apply}: Apply a function over an H2O parsed data object (an array).\newline
%%\newpage
\emph{Group By Aggregation}\par
{\texttt{h2o.ddply}}: Split H2O dataset, apply a function, and display results.\newline
{\texttt{h2o.group\_by}}: Apply an aggregate function to each group of an H2O dataset.\newline

\emph{Tabulation}\par
{\texttt{h2o.table}}: Use the cross-classifying factors to build a table of counts at each combination of factor levels.

\subsection{Data Munging}\par

\emph{General Column Manipulations}\par
{\texttt{is.na}}: Display missing elements.  \newline

\emph{Element Index Selection}\par 
{\texttt{h2o.which}}: Display the row numbers for which the condition is true.\newline

\emph{Conditional Element Value Selection}\par 
{\texttt{h2o.ifelse}}: Apply conditional statements to numeric vectors in H2O parsed data objects.\newline

\emph{Numeric Column Manipulations}\par
{\texttt{h2o.cut}}: Convert H2O Numeric Data to Factor. \newline

\emph{Character Column Manipulations}\par
{\texttt{h2o.strsplit}}: Splits the given factor column on the input split. \newline
{\texttt{h2o.tolower}}: Change the elements of a character vector to lower case. \newline
{\texttt{h2o.toupper}}: Change the elements of a character vector to lower case. \newline
{\texttt{h2o.trim}}: Remove leading and trailing white space.\newline
{\texttt{h2o.gsub}}: Match a pattern \& replace all instances of the matched pattern with the replacement string globally. \newline
{\texttt{h2o.sub}}: Match a pattern \& replace the first instance of the matched pattern with the replacement string.\newline

\emph{Factor Level Manipulations}\par
{\texttt{h2o.levels}}: Display a list of the unique values found in a column of categorical data. \newline

\emph{Date Manipulations}\par
{\texttt{h2o.month}}: Convert the entries of a H2OParsedData object from milliseconds to months (on a 0 to 11 scale). \newline
{\texttt{h2o.year}}: Convert the entries of a H2OParsedData object from milliseconds to years, indexed starting from 1900.\newline

\emph{Matrix Operations}\par
{\texttt{\%$*$\%}}}: Multiply two matrices, if they are conformable.\newline
{\texttt{t}}: Given a matrix or data.frame x, t returns the transpose of x.\newline

\subsection{Data Modeling}

\emph{Model Training: Supervised Learning}\par
{\texttt{h2o.deeplearning}}: Perform Deep Learning neural networks on an H2OParsedData object.\newline
{\texttt{h2o.gbm}}: Build gradient boosted classification trees and gradient boosted regression trees on a parsed data set. \newline
{\texttt{h2o.glm}}: Fit a generalized linear model, specified by a response variable, a set of predictors, and a description of the error distribution. \newline
{\texttt{h2o.naiveBayes}}: Build gradient boosted classification trees and gradient boosted regression trees on a parsed data set.\newline
{\texttt{h2o.prcomp}}: Perform principal components analysis on the given data set. \newline
{\texttt{h2o.randomForest}}: Perform random forest classification on a data set.\newline

\emph{Model Training: Unsupervised Learning}\par
{\texttt{h2o.anomaly}}: Detect anomalies in a H2O dataset using a H2O deep learning model with auto-encoding.\newline
{\texttt{h2o.deepfeatures}}: Extract the non-linear features from a H2O dataset using a H2O deep learning model.\newline
{\texttt{h2o.kmeans}}: Perform k-means clustering on a data set. \newline
%%\newpage
\emph{Grid Search}\par
{\texttt{h2o.grid}}: Efficient method to build multiple models with different hyperparameters. \newline

\emph{Model Scoring}\par
{\texttt{h2o.predict}}: Obtain predictions from various fitted H2O model objects.\newline

\emph{Classification Model Helpers}\par
{\texttt{h2o.accuracy}}: Get the between cluster sum of squares.\newline
{\texttt{h2o.auc}}: Retrieve the AUC (area under ROC curve).\newline
{\texttt{h2o.confusionMatrix}}: Display prediction errors for classification data from a column of predicted responses and a column of actual (reference) responses in H2O.\newline
{\texttt{h2o.hit\_ratio\_table}}: Retrieve the Hit Ratios. If {\texttt{train}, {\texttt{valid}}, and {\texttt{xval}} parameters are FALSE (default), then the training Hit Ratios value is returned. If more than one parameter is set to TRUE, then a named list
of Hit Ratio tables are returned, where the names are {\texttt{train}, {\texttt{valid}}, or {\texttt{xval}}.\newline
{\texttt{h2o.performance}}: Evaluate the predictive performance of a model via various measures.\newline

\emph{Regression Model Helper}\par
{\texttt{h2o.mse}}: Display the mean squared error calculated from a column of predicted responses and a column of actual (reference) responses in H2O.\newline

\emph{Clustering Model Helper}\par
{\texttt{h2o.betweenss}}: Get the between cluster sum of squares.\newline
{\texttt{h2o.centers}}: Retrieve the Model Centers.\newline

\subsection{H2O Cluster Operations}
\emph{H2O Key Value Store Access}\par
{\texttt{h2o.assign}}: Assign H2O hex.keys to objects in their R environment.\newline
{\texttt{h2o.getFrame}}: Get a reference to an existing H2O data set. \newline
{\texttt{h2o.getModel}}: Get a reference to an existing H2O model. \newline
{\texttt{h2o.ls: }}Display a list of object keys in the running instance of H2O. \newline
{\texttt{h2o.rm}}: Remove H2O objects from the server where the instance of H2O is running, but does not remove it from the R environment.\newline

\emph{H2O Object Serialization}\par
{\texttt{h2o.loadModel}}: Load an H2OModel object from disk.\newline
{\texttt{h2o.saveModel}}: Save an H2OModel object to disk to be loaded back into H2O using {\texttt{h2o.loadModel}}.\newline

\emph{H2O Cluster Connection}\par
{\texttt{h2o.init (nthreads = -1)}}: Connect to a running H2O instance using all CPUs on the host and check the local H2O R package is the correct version.\newline
{\texttt{h2o.shutdown}}: Shut down the specified H2O instance. All data on the server will be lost!\newline

\emph{H2O Load Balancing}\par
{\texttt{h2o.rebalance}}: Rebalance (repartition) an existing H2O data set into given number of chunks (per Vec), for load-balancing across multiple threads or nodes.\newline %%still supported? couldn't find in R doc

\emph{H2O Cluster Information}\par
{\texttt{h2o.clusterInfo}}: Display the name, version, uptime, total nodes, total memory, total cores and health of a cluster running H2O.\newline
{\texttt{h2o.clusterStatus}}: Retrieve information on the status of the cluster running H2O.\newline

\emph{H2O Logging}\par
{\texttt{h2o.clearLog}}: Clear all H2O R command and error response logs from the local disk. \newline
{\texttt{h2o.downloadAllLogs}}: Download all H2O log files to the local disk.\newline
{\texttt{h2o.logAndEcho}}: Write a message to the H2O Java log file and echo it back.  \newline
{\texttt{h2o.openLog}}: Open existing logs of H2O R POST commands and error responses on the local disk.\newline
{\texttt{h2o.getLogPath}}: Get the file path for the H2O R command and error response logs.\newline
{\texttt{h2o.startLogging}}: Begin logging H2O R POST commands and error responses. \newline
{\texttt{h2o.stopLogging}}: Stop logging H2O R POST commands and error responses.\newline

\emph{H2O String Manipulation}\par
{\texttt{h2o.gsub}}: String global substitution (all occurrences).\newline
{\texttt{h2o.strsplit}}: String Split.\newline
{\texttt{h2o.sub}}: String substitution (first occurrence).\newline
{\texttt{h2o.tolower}}: Convert characters to lower case.\newline
{\texttt{h2o.toupper}}: Convert characters to upper case.\newline
{\texttt{h2o.trim}}: Trim spaces.\newline


\end{document}